% Módulo científico sin preámbulo ni portada autónoma.
% Dependencia causal: atlas, operadores K_D/K_\Omega y contraste metrológico.

\chapter{Realización equivariante de especies y operadores de masa: fibras genealógicas, medidas espectrales, polos complejos y composiciones materiales}
\label{hmtmasscl:chap-realizacion-equiv}

\section{Dominio interno de la realización material y clasificación de sus codominios}

\subsection{Separación entre línea espectral, multiplete, distribución, polo complejo, objeto categórico y criterio de imposibilidad}

Sea \(\mathcal P\) el conjunto de estados físicos individualizados. Un estado
elemental se representa mediante el descriptor
\[
 X=(\mathfrak c,\rho,g,f,q,\epsilon,\varsigma,\mathfrak a),
\]
que conserva clase de representación \(\mathfrak c\), representación interna
\(\rho\), generación \(g\), sabor \(f\), carga \(q\), orientación de hoja
\(\epsilon\), sector estadístico \(\varsigma\) y datos de acoplamiento
\(\mathfrak a\). Un estado compuesto añade constituyentes ordenados, árbol de
composición, frontera común y números cuánticos colectivos. Una resonancia
añade los canales abiertos y la prescripción causal del resolvente.

La realización física debe responder a tres preguntas distintas:

\begin{enumerate}
\item qué fibra, órbita o multisección del atlas corresponde al descriptor;
\item qué operador positivo publica la masa sobre esa fibra;
\item qué acoplamiento produce un escalar, un multiplete, una distribución o
      un polo complejo.
\end{enumerate}

Confundir estas preguntas introduce un selector retrospectivo. La respuesta
correcta se obtiene mediante una relación natural, un operador de fibra y una
compresión espectral, en ese orden.

\subsection{Antecedencia del atlas genealógico respecto de toda identificación metrológica}

La Ley General de Masas no parte de una tabla de nombres ni de una lista de
valores experimentales. Su dominio interno es la cadena finita
\[
 \mathcal R_{324}\longrightarrow \mathcal C_{81}
 \longrightarrow \mathcal F_{56}\longrightarrow \mathcal M_{13},
\]
donde \(\mathcal R_{324}\) contiene las \(4\) orientaciones de cada una de las
\(81\) celdas, \(\mathcal F_{56}\) reúne historias de ruta equivalentes y
\(\mathcal M_{13}\) conserva las multisecciones operatorias. El sector nulo
se bifurca antes de aplicar el carácter positivo. Por tanto, ni una tabla
histórica abreviada ni una taxonomía de nombres físicos puede sustituir este
dominio.

\section{Producto fibrado entre descriptores de especie y rutas del atlas 81/56/13/324}

\subsection{Condición exacta de no vacuidad y selección prospectiva sin selector puntual ilícito}

Sea \(\mathcal D\) el espacio finito de datos internos ya construidos antes del
contraste: representación, orientación, memoria, frontera, carga, sabor,
color, generación y sector estadístico. Denotemos por
\[
 d_{\mathcal P}:\mathcal P\longrightarrow\mathcal D,
 \qquad
 d_{\mathcal R}:\mathcal R_{324}\longrightarrow\mathcal D
\]
los mapas de descriptor físico y descriptor genealógico. Se define el objeto
de realización por el producto fibrado
\[
 \mathcal J
 =\mathcal P\times_{\mathcal D}\mathcal R_{324}
 =\{(X,\gamma):d_{\mathcal P}(X)=d_{\mathcal R}(\gamma)\}.
\]
La fibra de una especie es
\[
\mathscr S(X)=\{\gamma\in\mathcal R_{324}:(X,\gamma)\in\mathcal J\}.
\]
Un descriptor (X) se denomina realizable por el atlas cuando
\[
 d_{\mathcal P}(X)\in\operatorname{im}d_{\mathcal R}.
\]
Esta condición es equivalente a \(\mathscr S(X)\ne\varnothing\). No se deduce
de la mera escritura del producto fibrado: constituye la condición exacta de
existencia que debe verificarse para cada descriptor físico.
El valor de \(\mathscr S\) puede ser una ruta, una órbita o una multisección.
No se fuerza una sección puntual cuando la simetría interna exige conservar
una multiplicidad.

\subsection{Equivarianza de las fibras, órbitas y multisecciones bajo refinamiento}

\begin{definition}[Relación de realización admisible]
La relación \(\mathscr S\) es admisible cuando satisface simultáneamente:
\begin{enumerate}
\item depende únicamente de datos construidos antes de abrir la tabla externa;
\item conmuta con la conjugación partícula--antipartícula;
\item es equivariante bajo la inversión celular y las acciones internas de
      calibre;
\item es compatible con refinamiento, retorno nonádico y lectura
      dodecafásica;
\item conserva toda multiplicidad no eliminada por un acoplamiento físico;
\item permanece invariante al permutar o reemplazar los valores externos de
      masa y anchura.
\end{enumerate}
\end{definition}

\begin{theorem}[Realización prospectiva por producto fibrado]
Si \(d_{\mathcal P}\) y \(d_{\mathcal R}\) son equivariantes respecto de una
simetría \(G\) y no dependen de magnitudes externas, entonces \(\mathcal J\)
es un \(G\)-subobjeto de \(\mathcal P\times\mathcal R_{324}\). Para cada
descriptor realizable \(X\), la fibra no vacía \(\mathscr S(X)\) es
\(G_X\)-estable y la asignación \(X\mapsto\mathscr S(X)\) es independiente del
resultado que posteriormente se contraste.
\end{theorem}

\begin{proof}
Para \((X,\gamma)\in\mathcal J\) se tiene
\(d_{\mathcal P}(X)=d_{\mathcal R}(\gamma)\). La equivariancia da
\[
 d_{\mathcal P}(gX)=g\,d_{\mathcal P}(X)
 =g\,d_{\mathcal R}(\gamma)=d_{\mathcal R}(g\gamma),
\]
y, por consiguiente, \((gX,g\gamma)\in\mathcal J\). Si \(g\in G_X\), la
segunda componente permanece en \(\mathscr S(X)\). La independencia de los
valores externos es inmediata porque ninguno de ellos pertenece al dominio
de los dos mapas del producto fibrado.
\end{proof}

\section{Clasificación espectral intrínseca de las especies materiales}

\subsection{Descomposición por soporte puntual, bloque irreducible y medida imagen de fibra}

Sea \(\mathcal A_{\rm rut}\) el álgebra conmutativa finita generada por los
observables internos de ruta: firma \(\mathbb Z^5\), carácter, cronología,
\(K_D\), \(K_\Omega\), orientación, memoria y lectores compatibles. En una
fibra \(F\), se define
\[
 \gamma\sim_{\rm sp}\eta
 \quad\Longleftrightarrow\quad
 a(\gamma)=a(\eta)\quad\text{para todo }a\in\mathcal A_{\rm rut}.
\]
Para cada clase \(\lambda\in F/\!\sim_{\rm sp}\),
\[
 P_\lambda=\sum_{\gamma\in\lambda}
 |\gamma\rangle\langle\gamma|
\]
es un proyector espectral; los \(P_\lambda\) son ortogonales y
\(\sum_\lambda P_\lambda=I_F\).

\begin{theorem}[Máxima individualización interna]
Todo clasificador construido exclusivamente con observables de
\(\mathcal A_{\rm rut}\) factoriza de manera única por el cociente
\(F/\!\sim_{\rm sp}\). En consecuencia, dos descriptores con firma completa
distinta pertenecen a sectores distinguibles; dos descriptores que coinciden
para todos los generadores del álgebra no pueden separarse mediante una regla
construida con esos mismos observables.
\end{theorem}

\begin{proof}
Los caracteres de la álgebra conmutativa finita separan exactamente su
espectro conjunto. Un clasificador funcional de sus generadores es constante
en cada fibra del mapa espectral conjunto y, por ello, factoriza por el
cociente. Si dos rutas tienen el mismo carácter sobre toda la álgebra, ningún
elemento de ella puede separarlas.
\end{proof}

\subsection{Distinción entre masa escalar, multiplete y distribución espectral}

Este teorema cierra el límite lógico de la clasificación. Cuando una
representación física exige distinguir dos estados pertenecientes a una misma
clase espectral, debe exhibir un observable adicional o un acoplamiento que
amplíe \(\mathcal A_{\rm rut}\). La ausencia de tal observable no autoriza a
usar la masa externa como etiqueta selectora.

\section{Refinamiento natural del registro enriquecido en las torres positivas \texorpdfstring{\(90/120\)}{90/120}}

\subsection{Enumeración positiva del semigrupo \texorpdfstring{\(\langle90,120\rangle\)}{<90,120>} y exclusión del modo 0}

La existencia del semigrupo \(\langle90,120\rangle\) no basta por sí sola:
debe construirse el mapa que asigna coeficientes a una ruta enriquecida. Sea
\[
 \mathcal H_{90,120}
 =\{90a+120b:a,b\in\mathbb N_0,\ (a,b)\ne(0,0)\}
 =\{90,120\}\cup\{30r:r\geq6\}
 =\{h_0<h_1<h_2<\cdots\}
\]
su enumeración estrictamente creciente, donde \(\mathbb N_0=\{0,1,2,\ldots\}\),
y sea \(s_{h_n}>0\) la normalización declarada del modo correspondiente. El
registro compatible de una ruta
\(\Gamma\) proporciona una sucesión de datos enteros finitos
\[
 \mathcal L(\Gamma)=(\ell_0(\Gamma),\ell_1(\Gamma),\ldots),
\]
en la que cada \(\ell_n\) reúne retorno, acarreo, devanado, memoria y frontera
en el refinamiento \(n\). Para fijar sin ambigüedad la codificación, sea
\[
 \nu(z)=
 \begin{cases}2z,&z\geq0,\\-2z-1,&z<0,\end{cases}
 \qquad
 \pi_{\rm C}(u,v)=\frac{(u+v)(u+v+1)}2+v,
\]
y defínase \(\iota:\mathbb Z^d\to\mathbb N_0\) por aplicación recursiva de
\(\pi_{\rm C}\) a \(\nu(z_1),\ldots,\nu(z_d)\). Entonces
\[
 q_n(\Gamma)
 =\frac{\iota(\ell_n(\Gamma))}
 {1+|\iota(\ell_n(\Gamma))|}\in(-1,1).
\]
La realización mínima de torre es
\[
 \eta_{h_n}(\Gamma)
 =\frac{3^{-(n+1)}q_n(\Gamma)}{s_{h_n}},
 \qquad
 \mathfrak T(\Gamma)
 =(\eta_h(\Gamma))_{h\in\mathcal H_{90,120}}.
\]

\subsection{Convergencia absoluta, naturalidad por prefijos y cota efectiva de cola}

\begin{theorem}[Convergencia, naturalidad y ceguera metrológica]
Para toda ruta enriquecida,
\[
 \sum_{h\in\mathcal H_{90,120}}|\eta_h(\Gamma)s_h|
 \leq\sum_{n\ge0}3^{-(n+1)}=\frac12.
\]
Si el refinamiento prolonga el registro por prefijos, \(\mathfrak T\) es natural
respecto de las truncaciones. Además, ningún valor de masa externo interviene
en \(\mathfrak T\), y la cola después de \(N\) modos satisface
\[
 \sum_{n>N}|\eta_{h_n}s_{h_n}|
 \leq\frac{3^{-(N+1)}}2.
\]
\end{theorem}

\begin{proof}
Se tiene \(|q_n|<1\); las dos cotas son sumas geométricas. La compatibilidad
por prefijos conserva todos los coeficientes ya construidos y añade los del
nuevo nivel, lo cual da el cuadrado de naturalidad con la truncación de
\(\mathcal E_{90,120}\). La fórmula sólo contiene el registro, la codificación
explícita y la normalización modal declarada.
\end{proof}

La construcción anterior fija una realización mínima explícita y convergente,
relativa a la codificación \(\iota\) y a las normalizaciones \(s_h\) declaradas.
Una dinámica que modifique sus pesos debe declarar el operador que efectúa esa
modificación y demostrar de nuevo naturalidad y convergencia; no puede
reconstruir los pesos a partir de un residual experimental.

\section{Operadores bidireccionales \texorpdfstring{\(K_D\) y \(K_\Omega\)}{K_D y K_Omega} y medida espectral conjunta}

\subsection{PVM conjunta, compresión a POVM y distribución escalar inducida por un estado}

Para \(F_X=\mathscr S(X)\), sea
\[
 \mathcal H_X=\ell^2(F_X)\otimes V_{\rho_X}
\]
el espacio de realización, donde \(V_{\rho_X}\) porta la representación
interna correspondiente. Cada ruta \(\gamma\in F_X\) posee firma
\[
 \sigma(\gamma)
 =(n_A,s_C,k_{\Delta_4},b_{120},b_\zeta)\in\mathbb Z^5
\]
y coeficientes de refinamiento
\(\eta(\gamma)\in\mathcal E_{90,120}\). El operador basal es
\[
 \widehat M_X^{(0)}
 =\sum_{\gamma\in F_X}
 m_{\rm clk}\,
 \chi_{\rm full}(\sigma(\gamma),\eta(\gamma))
 |\gamma\rangle\langle\gamma|\otimes I_{V_{\rho_X}}.
\]
La escala \(m_{\rm clk}\) desciende de la rama absoluta
acción--reloj--masa, incluida la década \(10^{-34}\); no se extrae de una
masa experimental.

Los lectores \(K_D\) y \(K_\Omega\) inducen operadores diagonales exactos
\(B_X^D\) y \(B_X^\Omega\). Las dos orientaciones publican
\[
 \widehat M_X^{\beta,r}
 =R_{\rm act}^{B_X^r/2}\widehat M_X^{(0)}
  R_{\rm act}^{B_X^r/2},
 \qquad r\in\{D,\Omega\}.
\]
Ambas son positivas y conmutan en la base de rutas. Por el teorema espectral
conjunto existe una única medida espectral proyectiva
\(E_X^{D,\Omega}\) sobre \(\mathbb R_+^2\) tal que
\[
 \widehat M_X^{\beta,D}
 =\int x\,dE_X^{D,\Omega}(x,y),
 \qquad
 \widehat M_X^{\beta,\Omega}
 =\int y\,dE_X^{D,\Omega}(x,y).
\]
En la base de rutas,
\[
 E_X^{D,\Omega}(A)|\gamma\rangle
 =\mathbf 1_A\!\left(m_D(\gamma),m_\Omega(\gamma)\right)
  |\gamma\rangle.
\]
Ésta es la salida bidireccional canónica: conserva el par completo y no
supone de antemano un funcional escalar.

Para un canal \(P_{X,a}\), la medida espectral comprimida
\[
 F_{X,a}^{D,\Omega}(A)
 =\left.P_{X,a}E_X^{D,\Omega}(A)P_{X,a}\right|_{P_{X,a}\mathcal H_X}
\]
es una POVM sobre el subespacio preparado. Para una densidad \(\rho_X\)
soportada en ese subespacio, la distribución escalar conjunta es
\[
 \mu_{X,a}^{D,\Omega}(A)
 =\operatorname{Tr}\!\left(
 \rho_X F_{X,a}^{D,\Omega}(A)
 \right).
\]
Todo funcional declarado \(f:\mathbb R_+^2\to\mathbb R_+\) induce
\(\mu_{X,a}^f=f_\#\mu_{X,a}^{D,\Omega}\) sin modificar la PVM de origen.

\subsection{Momentos mixtos, soporte espectral puntual y conservación de multiplicidades}

Sea
\[
 \boldsymbol\beta_X(\gamma)
 =(m_D(\gamma),m_\Omega(\gamma)).
\]
La medida imagen uniforme de la fibra es
\[
 \nu_X=(\boldsymbol\beta_X)_\#
 \left(\frac1{|F_X|}\sum_{\gamma\in F_X}\delta_\gamma\right)
 =\frac1{|F_X|}\sum_{\gamma\in F_X}
 \delta_{(m_D(\gamma),m_\Omega(\gamma))}.
\]
Su transformada y sus momentos mixtos son
\[
 Z_X(s,t)=\int e^{sx+ty}\,d\nu_X(x,y),
 \qquad
 \mathfrak m_{p,q}(X)=\int x^py^q\,d\nu_X(x,y).
\]
Como \(\nu_X\) tiene soporte finito, \(Z_X\), o equivalentemente la familia
completa de momentos mixtos, determina el multiconjunto conjunto de pares
lectores hasta permutación de rutas. La medida imagen es, por tanto, un
invariante completo del par operatorial; no es una media de masas.

Si \(u:F_X\to F_Y\) es una biyección que conserva registro, orientación y
ambos lectores, el unitario \(U_u|\gamma\rangle=|u\gamma\rangle\) satisface
\[
 U_uE_X^{D,\Omega}(A)=E_Y^{D,\Omega}(A)U_u,
 \qquad
 \nu_X=\nu_Y.
\]
La misma igualdad vale para las distribuciones físicas cuando densidad y
proyector se transportan por \(U_u\). Esta naturalidad se afirma sólo para
los entrelazadores efectivamente construidos; una coincidencia de
cardinalidades no crea por sí sola una equivalencia de fibras.

\section{Escalarización por simetría y descomposición isótípica de las fibras}

\subsection{Esperanza condicional de simetría y criterio de reducción a una línea espectral}

Sea \(G_X\) el grupo compacto de simetría no rota de la fibra y
\(U_X:G_X\to\mathcal U(\mathcal H_X)\) su representación. La esperanza
condicional de simetría se define por
\[
 \mathbb E_X(T)
 =\int_{G_X}U_X(g)TU_X(g)^*\,d\mu_X(g),
\]
donde \(\mu_X\) es la medida de Haar normalizada. El operador físico de masa
antes de abrir canales se define, cuando el tipo de estado declara un
funcional \(\Phi_X:\mathbb R_+^2\to\mathbb R_+\), por
\[
 \mathcal M_X
 =\mathbb E_X\!\left(
   \Phi_X(\widehat M_X^{\beta,D},\widehat M_X^{\beta,\Omega})
  \right).
\]
El cálculo funcional se realiza con la medida conjunta
\(E_X^{D,\Omega}\). En una representación que contiene un entrelazador de
intercambio de lectores puede tomarse \(\Phi_X(x,y)=\sqrt{xy}\). Si el tipo
físico no declara \(\Phi_X\), la salida completa sigue siendo la PVM conjunta;
no se inventa un escalar.

\begin{theorem}[Cierre escalar, multiplete y distribución]
El operador \(\mathcal M_X\) pertenece al conmutante
\(U_X(G_X)'\). Si
\[
 \mathcal H_X\simeq
 \bigoplus_\lambda V_\lambda\otimes\mathbb C^{m_\lambda},
\]
entonces
\[
 \mathcal M_X=
 \bigoplus_\lambda I_{V_\lambda}\otimes M_{X,\lambda}.
\]
En un sector irreducible de multiplicidad \(1\), la restricción es un escalar.
En un sector reducible, la salida exacta es el espectro de los bloques
\(M_{X,\lambda}\), con sus multiplicidades. No existe una indeterminación:
existe un codominio distinto.
\end{theorem}

\begin{proof}
La invariancia de la medida de Haar implica
\(U_X(h)\mathbb E_X(T)U_X(h)^*=\mathbb E_X(T)\) para todo \(h\in G_X\).
La descomposición sigue del teorema de completa reducibilidad y del lema de
Schur. El caso irreducible produce un múltiplo de la identidad; el caso
reducible conserva exactamente los bloques de multiplicidad.
\end{proof}

Sea \(P_{X,a}\) el proyector determinado por el estado preparado y el canal o
aparato físico \(a\). La medida espectral observable es
\[
 \mu_{X,a}(B)=
 \langle\Psi_X,
 P_{X,a}E_{\mathcal M_X}(B)P_{X,a}\Psi_X\rangle.
\]
Hay una línea escalar precisamente cuando el soporte espectral comprimido es
puntual:
\[
 E_{\mathcal M_X}(\{m_X\})P_{X,a}=P_{X,a},
 \qquad\text{equivalentemente}\qquad
 (\mathcal M_X-m_X I)\,P_{X,a}=0.
\]
La igualdad aislada
\(P_{X,a}\mathcal M_X P_{X,a}=m_X P_{X,a}\) sólo fija una esperanza si el
subespacio no es reductor y no basta para demostrar una línea espectral.
Si la compresión conserva varios autovalores, la salida es un multiplete o una
distribución. Sustituirla por una media numérica destruiría información
física y genealógica.

\subsection{Media geométrica de los lectores únicamente bajo invariancia de intercambio demostrada}

Cuando la representación física exige simetría bajo el intercambio
\(D\leftrightarrow\Omega\), el agregado operatorial mínimo es la media
geométrica
\[
 \widehat M_X^{\rm sym}
 =\widehat M_X^{\beta,D}\#\widehat M_X^{\beta,\Omega},
 \qquad
 A\#B
 =A^{1/2}\bigl(A^{-1/2}BA^{-1/2}\bigr)^{1/2}A^{1/2},
\]
sobre el soporte positivo común. En el caso diagonal,
\[
 \widehat M_X^{\rm sym}
 =R_{\rm act}^{(B_X^D+B_X^\Omega)/4}
  \widehat M_X^{(0)}
  R_{\rm act}^{(B_X^D+B_X^\Omega)/4}.
\]
Esta fórmula es consecuencia de la hipótesis explícita de intercambio; no
sustituye la medida conjunta cuando tal simetría no ha sido establecida.

\begin{theorem}[Caracterización del agregado simétrico]
La media geométrica es positiva, simétrica, covariante por congruencia para
toda \(S\) invertible y autodual:
\[
 A\#B=B\#A,
 \qquad
 S^*(A\#B)S=(S^*AS)\#(S^*BS),
 \qquad
 (A\#B)^{-1}=A^{-1}\#B^{-1}.
\]
Además, \(A\#A=A\). Por consiguiente, en los sectores donde existe un
entrelazador físico que intercambia los lectores, el agregado no introduce un
parámetro ni privilegia una orientación.
\end{theorem}

\begin{proof}
La identidad de congruencia se entiende para \(S\) invertible. Las identidades
resultan del cálculo funcional de la raíz positiva de
\(A^{-1/2}BA^{-1/2}\). En el caso conmutativo se reducen a
\(A\#B=(AB)^{1/2}\).
\end{proof}

Existe además un agregado escalar canónico bajo axiomas distintos. Si
\(s_n\) es continuo, invariante por permutación, multiplicativo coordenada a
coordenada y normalizado por \(s_n(cI)=c\), entonces
\[
 s_n(T)=\det(T)^{1/n}.
\]
En efecto, al pasar a logaritmos, continuidad y aditividad hacen lineal el
funcional sobre los logaritmos de los autovalores; la simetría iguala todos
los coeficientes y la normalización fija su valor en \(1/n\). Para el par de
lectores, la simetría de intercambio da el invariante
\[
 S_{D,\Omega}
 =\sqrt{
 \operatorname{ndet}(T_D)\operatorname{ndet}(T_\Omega)},
 \qquad
 \operatorname{ndet}(T)=\det(T)^{1/n}.
\]
Este resultado clasifica la fibra completa bajo los axiomas enunciados. No se
identifica automáticamente con una línea de masa: dicha identificación exige
que el acoplamiento físico adopte precisamente ese funcional.

\section{Sector nulo y antecedencia de la carta de masa cero respecto del carácter positivo}

\subsection{Separación de las representaciones fotónica \texorpdfstring{\(U(1)\)}{U(1)} y gluónica adjunta \texorpdfstring{\(SU(3)\)}{SU(3)}}

El fotón pertenece a la carta nula electromagnética y los gluones a la carta
nula adjunta de color. Sus diferentes representaciones, polarizaciones y
restricciones de calibre permanecen visibles aunque ambas masas sean nulas.

\subsection{No alcanzabilidad del sector nulo como límite de la rama positiva de masa}

La nulidad no se obtiene tomando un límite del carácter positivo. El dominio
se separa primero en
\[
 \mathcal H_{\rm phys}=\mathcal H_0\oplus\mathcal H_+,
\]
con
\[
 \mathcal M|_{\mathcal H_0}=0,
 \qquad
 \mathcal M|_{\mathcal H_+}>0.
\]

\section{Clausura beta del electrón y sectores leptónicos cargados}

\subsection{Unicidad del electrón en la celda \texorpdfstring{\((5,5)\)}{(5,5)} y preservación de la masa beta rectora}

El electrón constituye la excepción central. La inversión celular posee un
único punto fijo, la celda \((5,5)\), y por ello su fibra física es de rango
uno. Las restantes multisecciones estables no admiten, en general, un selector
puntual equivariante. Esta diferencia no es una carencia del atlas: es la
distinción matemática entre una sección central y una especie realizada como
órbita o multisección.

Para la fibra electrónica de rango uno, la medida conjunta es puntual, los dos
lectores comparten el primer componente \(80\), la compresión es escalar y se
recupera el cierre beta
\[
 m_e^\beta
 =0.5109989506900008086082571648\ldots\ {\rm MeV}.
\]
El valor basal permanece como etapa de construcción; no reemplaza al cierre
beta.

\subsection{Realizaciones del muón y del tau mediante sectores cargados y criterio espectral de escalaridad}

En el sector leptónico cargado, los descriptores distinguen la fibra central,
la multisección de profundidad \(G_2\) y la multisección de profundidad
\(G_4\). La primera produce la línea electrónica única; las otras dos
producen los operadores asociados a las realizaciones muónica y tauónica. Si
el acoplamiento preserva más de un bloque, la predicción correcta es su
espectro; la etiqueta física no autoriza a escoger uno de sus autovalores por
proximidad numérica.

\section{Álgebra cromática \texorpdfstring{\(SU(3)\)}{SU(3)}, sabores quark y transporte de la masa corriente}

\subsection{Álgebra qutrit de Weyl, reducción de Reynolds y proyectores singlete}

La órbita cromática se realiza en el qutrit de color. Los operadores de Weyl
generan \(\mathfrak{sl}_3(\mathbb C)\); la selección antihermítica de traza
nula produce \(\mathfrak{su}(3)\), y la conexión finita transporta las tres
cartas cromáticas. La masa de sabor debe ser invariante bajo esta acción.

Sea \(U_{\rm col}:SU(3)\to\mathcal U(V_{\rm col})\). La reducción cromática
es la esperanza de Reynolds
\[
 \mathbb E_{\rm col}(T)
 =\int_{SU(3)}U_{\rm col}(g)T U_{\rm col}(g)^*\,dg.
\]
En la representación fundamental irreducible,
\[
 \mathbb E_{\rm col}(\mathcal M_q)
 =m_q(\mu,\mathscr R)I_3.
\]
La masa publicada es, por tanto, independiente de rojo, verde y azul. Las
coordenadas \(\mu\) y \(\mathscr R\) identifican escala y esquema de
renormalización; no son ambigüedades ocultas y deben acompañar siempre a la
cifra.

Para mesones y bariones, los proyectores de singlete son
\[
 P_{q\bar q}^{(1)}
 =\frac1{3}|\Omega\rangle\langle\Omega|,
 \qquad
 |\Omega\rangle=\sum_{a=1}^{3}|a\bar a\rangle,
\]
y
\[
 P_{qqq}^{(1)}
 =|\varepsilon\rangle\langle\varepsilon|,
 \qquad
 |\varepsilon\rangle
 =\frac1{\sqrt6}\sum_{a,b,c}\varepsilon_{abc}|abc\rangle.
\]
Ambos eliminan la etiqueta cromática observable sin borrar la órbita interna
que intervino en la construcción.

\subsection{Conexión de renormalización y transporte ordenado entre fibras de escala}

La dependencia en escala es una sección de un fibrado sobre
\(t=\log\mu\). Si \(\gamma_m(g(t))\) es el endomorfismo inducido por la
conexión cromática, el transporte se define por
\[
 \nabla_{\rm RG}m_q
 =\left(\frac{d}{dt}+\gamma_m(g(t))\right)m_q=0,
\]
\[
 m_q(t_2)
 =\mathcal P\exp\!\left(-\int_{t_1}^{t_2}
 \gamma_m(g(t))\,dt\right)m_q(t_1).
\]
Dos valores pertenecientes a escalas o esquemas distintos no son números
comparables hasta transportarlos a la misma fibra. Esta ecuación cierra el
tipo matemático de la masa corriente: no es un escalar absoluto aislado, sino
una sección covariantemente transportada.

\section{Mezcla CKM y covariancia de los marcos espectrales quark}

\subsection{Transporte unitario del operador \texorpdfstring{\(B_d\)}{B_d} en el marco espectral de tipo arriba mediante \texorpdfstring{\(V_{\rm CKM}\)}{V_CKM}}

La mezcla CKM pertenece a un mapa distinto. Si \(B_u\) y \(B_d\) son los
operadores de masa en las bases de tipo arriba y tipo abajo, la matriz unitaria
interna satisface
\[
 B_d^{(u)}=V_{\rm CKM} B_d V_{\rm CKM}^*.
\]

\subsection{Separación entre el cierre interno y el reconocimiento estadístico exterior}

Los ángulos derivados de \(A\), \(C^*\) y \(\Delta_4\), la fase
\(\delta_{\rm CP}\), el invariante de Jarlskog y la identidad del determinante
del conmutador determinan la mezcla sin crear los autovalores de masa. Así se
mantienen separados espectro, color, generación y cambio de base.

\section{Operador neutrínico, mezcla PMNS y criterios Dirac--Majorana}

\subsection{Espectro neutral tridimensional y herencia de la escala acción--reloj--masa}

Los neutrinos son secciones fermiónicas eléctricamente neutras. La
fermionicidad desciende de la holonomía espinorial, la completación exterior y
las relaciones canónicas de anticonmutación; la neutralidad no altera su
estadística. El soporte interno es
\[
 \mathcal H_\nu=
 \bigoplus_{j=1}^{4}\mathcal H_{\nu,G_j},
 \qquad
 \dim\mathcal H_\nu=2+4+6+8=20.
\]
Las cuatro profundidades estructurales no son sabores ni autovalores. El
subespacio físico tridimensional se obtiene mediante el proyector neutral
\(P_\nu\), y el operador de masa es
\[
 M_\nu=P_\nu\mathcal M_\nu P_\nu.
\]
Sus tres autovalores ordenados son las masas; la escala absoluta procede del
mismo reloj másico que fija las restantes ramas, y no de los ángulos de mezcla.

\subsection{Rango 3 de la matriz de Gram, marco polar y criterio antiunitario de Majorana}

Sea \(M_\ell=P_\ell\mathcal M_\ell P_\ell\) el operador cargado. Denotemos por
\(F_\ell\) y \(F_\nu\) marcos espectrales ortonormales completos, ordenados por
los proyectores espectrales de \(M_\ell\) y \(M_\nu\). La matriz de mezcla es
\[
 U_{\rm PMNS}=F_\ell^*F_\nu.
\]

\begin{theorem}[Cierre espectral de la mezcla leptónica]
Si los dos marcos son completos, \(U_{\rm PMNS}\) es unitaria. Si los
espectros son simples, queda determinada salvo reajustes diagonales de fase y
permutaciones fijadas por el orden espectral. Si existe degeneración, el
resultado exacto es una clase doble
\[
 \prod_\lambda U(m_{\ell,\lambda})
 \backslash U(3) /
 \prod_\mu U(m_{\nu,\mu}),
\]
y no una matriz numérica artificialmente seleccionada.
\end{theorem}

\begin{proof}
La ortonormalidad da
\(U_{\rm PMNS}^*U_{\rm PMNS}=F_\nu^*F_\ell F_\ell^*F_\nu=I_3\).
Los cambios de base dentro de cada subespacio degenerado actúan por la
izquierda o por la derecha, lo que produce exactamente la clase doble
indicada.
\end{proof}

El registro completo determina los elementos de \(M_\nu\) mediante acarreo,
hoja, orientación, vacancia y holonomía. De modo equivalente, puede emplearse
el núcleo hermítico
\[
 K_{\rm reg}(c,d)
 =\sum_{r\in\{\rm car,hoj,ori,vac,hol\}}
 \frac{\langle\Xi_r(c),\Xi_r(d)\rangle}
      {\sum_x\|\Xi_r(x)\|^2},
\]
con omisión de todo sumando de norma nula. Su matriz de Gram es positiva y
construye el operador antes de la diagonalización. La fase de rango uno
\[
 I+(e^{i\pi/6}-1)|b_K\rangle\langle b_K|
\]
permanece como control negativo: su unitariedad formal no basta y su ángulo
\(\theta_{13}\) falsifica su identificación con la mezcla física de rango
completo.

Sea \(G_0\) la matriz tridimensional obtenida al agregar el núcleo anterior
sobre los tres sectores neutrínicos. La condición finita
\[
 \det G_0\ne0
\]
es el certificado exacto de que los observables incluidos en el registro
separan tres direcciones. En tal caso, el factor polar
\[
 U_{G_0}=G_0(G_0^*G_0)^{-1/2}
\]
es unitario y de rango \(3\). Si \(\det G_0=0\), el núcleo declarado no
determina un transporte tridimensional: debe añadirse un observable interno
independiente y volver a calcular la matriz de Gram. No es lícito imponer la
invertibilidad añadiendo un término construido a partir de la matriz de mezcla
que se pretende obtener. La matriz física queda fijada por los marcos
espectrales de \(M_\ell\) y \(M_\nu\); el criterio de Gram certifica suficiencia
del registro, no selecciona por sí solo esos marcos.

Sea \(\mathcal C_\nu\) la autorconjugación antiunitaria del sector neutral. El
criterio de Majorana es
\[
 \mathcal C_\nu^2=I,
 \qquad
 \mathcal C_\nu M_\nu=M_\nu\mathcal C_\nu,
 \qquad
 \mathcal C_\nu H_\nu=H_\nu\mathcal C_\nu,
\]
junto con la existencia de una base real fija en la realificación del
subespacio físico. Cuando se cumplen estas condiciones, cada modo puede
representarse por \(\psi=\mathcal C_\nu\psi\). La rama de Dirac exige, en
cambio, una carga conservada \(Q\) tal que
\[
 \mathcal C_\nu Q\mathcal C_\nu^{-1}=-Q,
 \qquad Q\psi\ne0,
\]
y conserva como modos distintos \(\psi\) y \(\mathcal C_\nu\psi\). Si no se
acredita ninguno de ambos conjuntos de condiciones, el tipo no queda decidido.
La neutralidad, la involución de palabra o el sector de Ising, tomados por
separado, no deciden esta alternativa.

\section{Sector neutrínico estéril y separación entre existencia y mezcla}

\subsection{Proyector ortogonal al sector activo y anulación de la corriente débil}

Una profundidad estructural adicional no constituye por sí sola un neutrino
estéril. Un sector estéril candidato queda definido por un proyector no nulo
\(P_s\) que satisface
\[
 P_s P_{\rm act}=0,
 \qquad P_s J_{\text{débil}}=0,
\]
con dinámica definida sobre \(P_s\mathcal H_\nu\).

\subsection{Distinción entre reducción invariante y acoplamiento no nulo con el sector activo}

Si además \([P_s,M_\nu]=0\), el sector es reductor y no mezcla con el sector activo. Si
\(P_sM_\nu P_{\rm act}\ne0\), existe mezcla activa--estéril y necesariamente
\([P_s,M_\nu]\ne0\). Por tanto, la existencia del sector y la existencia de
mezcla son propiedades distintas. Si no existe el proyector débilmente neutro,
el canal grueso adicional es profundidad interna del soporte y no una cuarta
especie.

\section{Bosones \texorpdfstring{\(W\) y \(Z\)}{W y Z} y sector escalar: proyectores, resolventes y criterio de polo}

\subsection{Operador efectivo sectorial y continuación meromorfa del resolvente reducido}

Los sectores \(W\), \(Z\) y escalar poseen descriptores propios de
representación, ruta y firma. Su realización no exige forzarlos dentro de una
fibra elemental de rango uno. Se construye primero el operador positivo de
fibra y después, en un subespacio \(P_X\mathcal H_X\) de dimensión finita, el
operador efectivo de canales
\[
 H_{X,\rm eff}(z)
 =P_X H_X P_X
 +P_X H_X Q_X(z-Q_X H_X Q_X)^{-1}Q_X H_X P_X.
\]
\subsection{Cero aislado admisible, residuo físico y distinción respecto de la mera apertura de canal}

Si el resolvente reducido admite continuación meromorfa a través del corte del
continuo, los polos son los ceros aislados, contados con multiplicidad, de
\[
 \det\bigl(z-H_{X,\rm eff}(z)\bigr)=0,
 \qquad
 z_X=M_X c_{\rm int}^2-\frac{i}{2}\Gamma_X.
\]
Un canal abierto no garantiza por sí solo la existencia de tales ceros. Cuando
existe un cero con el residuo físico prescrito, la parte real publica la masa de
polo y la parte imaginaria la anchura. Para el
sector escalar, la identificación de una excitación exige además el proyector
escalar y sus acoplamientos de calibre y fermiónicos. Dos polos independientes
determinan dos estados; una segunda etiqueta sin proyector ni polo no crea una
partícula.

\section{Estados compuestos y ley cocíclica de enlace}

\subsection{Proyectores singlete mesónicos y bariónicos y refinamiento de la firma compuesta}

Para constituyentes \(X_1,\ldots,X_n\), la entrada es el registro compuesto
\[
 \mathfrak C_{12}^{(n)}:
 (\Gamma_1,\ldots,\Gamma_n;\mathcal T,\partial)
 \longmapsto\mathcal L_{\rm comp},
\]
no la suma de masas aisladas. La composición binaria
\[
 L_1\star L_2=L_1+L_2+\mathfrak b(L_1,L_2)
\]
es asociativa exactamente cuando
\[
 \mathfrak b(L_1,L_2)+\mathfrak b(L_1\star L_2,L_3)
 =\mathfrak b(L_2,L_3)+\mathfrak b(L_1,L_2\star L_3).
\]
La firma, el refinamiento y el operador compuesto se calculan después del
enlace:
\[
 \mathcal L_{\rm comp}\longmapsto\sigma_{\rm comp}
 \longmapsto\eta_{\rm comp}\longmapsto\mathcal M_{\rm comp}.
\]
El proyector físico reúne singlete de color, espín y paridad, isospín, carga,
simetría de intercambio y canal.

\subsection{Reducción de Feshbach--Schur y construcción del operador efectivo}

La masa estable es un autovalor de la
compresión; cuando se cumplen las hipótesis de continuación meromorfa, la masa
y la anchura de una resonancia proceden de un polo de Feshbach. Con ello
mesones, bariones, estados exóticos, núcleos, átomos y
moléculas poseen el mismo esquema causal sin reducirse a una suma de masas.

Dos árboles de constituyentes representan la misma especie si existe un
unitario que entrelaza simultáneamente sus operadores, proyectores de canal y
polos. En ausencia de tal unitario se conservan como realizaciones distintas,
aunque compartan constituyentes y números cuánticos parciales.

\section{Resonancias, anchuras y supervivencia de las secciones materiales}

\subsection{Polos complejos, residuos admisibles y anchuras asociadas a la continuación meromorfa}

Sea \(S_X\) el proyector sobre el subespacio superviviente y
\(T_X=S_X U S_X\) el operador comprimido de un paso. Si posee un autovalor
dominante simple
\[
 \lambda_X=r_X e^{-i\theta_X},\qquad 0<r_X<1,
\]
fíjese la rama \(\theta_X\in I_\theta\), donde \(I_\theta\) es el intervalo
fundamental de cuasienergía determinado por la carta temporal y su orientación.
Entonces
\[
 M_X=\frac{\hbar_{\rm int}}{t_0c_{\rm int}^2}\theta_X,
 \qquad
 \Gamma_X=-\frac{2\hbar_{\rm int}}{t_0}\log r_X,
 \qquad
 \tau_X=\frac{\hbar_{\rm int}}{\Gamma_X}.
\]
\subsection{Compatibilidad del refinamiento genealógico con la persistencia resonante}

Sin la elección de \(I_\theta\), la fase sólo determina \(\theta_X\) módulo
\(2\pi\) y la masa de cuasienergía no es un escalar bien definido.
La anchura no es una sexta coordenada de la firma másica: procede del módulo
del modo abierto. La igualdad de masa y anchura entre partícula y
antipartícula exige transportar también canales y prescripción causal, no sólo
la paridad de la firma.

\section{Completación multiparticular y estadística de Fock}

\subsection{Funtores de Fock simétrico y exterior inducidos por la paridad de holonomía}

Sea \(\mathsf{Hol}_{\pm}^{\leq1}\) la categoría monoidal de fibras de una
partícula, con morfismos de refinamiento contractivos, en particular unitarios,
y paridad de holonomía
\(\epsilon\in\{+1,-1\}\). La completación de Fock se define por el funtor
\[
 \mathfrak F(H,\epsilon)=
 \begin{cases}
 \displaystyle\bigoplus_{n\ge0}\operatorname{Sym}^n(H),
     &\epsilon=+1,\\[2mm]
 \displaystyle\bigoplus_{n\ge0}\bigwedge^n\!H,
     &\epsilon=-1,
 \end{cases}
\]
y, para un entrelazador contractivo \(T:H\to H'\),
\[
 \mathfrak F(T)=
 \bigoplus_{n\ge0}\operatorname{Sym}^n(T)
 \quad\text{o}\quad
 \bigoplus_{n\ge0}\bigwedge^n\!T.
\]

\subsection{Relaciones CCR/CAR, exclusión de Pauli y dominio de los morfismos no contractivos}

\begin{theorem}[Transporte interno de la estadística]
La completación anterior es un funtor de operadores acotados sobre
\(\mathsf{Hol}_{\pm}^{\leq1}\), compatible con refinamiento y
monoidal graduada. En la rama \(\epsilon=-1\) los operadores de creación y
aniquilación satisfacen CAR y los operadores de ocupación cumplen
\(N_s^2=N_s\); en la rama \(\epsilon=+1\) satisfacen CCR. Por tanto, la
exclusión de Pauli y las estadísticas fermiónica y bosónica descienden de la
paridad de holonomía una vez fijada la representación de una partícula.
\end{theorem}

\begin{proof}
Las potencias simétricas y exteriores preservan composiciones e identidades;
la contractividad de \(T\) garantiza que su suma directa define un operador
acotado. Para morfismos no contractivos, la segunda cuantización requiere un
dominio denso explícito y no pertenece a esta categoría de operadores acotados.
La regla de signos del producto exterior produce CAR y la idempotencia de la
ocupación; el cociente simétrico produce CCR. La segunda cuantización de los
morfismos de refinamiento entrelaza los correspondientes operadores.
\end{proof}

Esta demostración cierra el transporte algebraico interno. La identificación
con el teorema relativista de espín--estadística requiere, además, el puente
lorentziano y la localidad declarados en Mecánica Dimensional; no deben
confundirse ambos enunciados.

\section{Realización termodinámica mediante estados de Gibbs y límite de presión}

\subsection{Hamiltoniano finito, partición, estado de Gibbs y límite cofinal de la presión}

Para un Hamiltoniano comprimido \(H_{X,V}\) sobre un volumen o refinamiento
finito \(V\), el estado térmico está completamente definido por
\[
 Z_{X,V}(\beta)=\operatorname{Tr}e^{-\beta H_{X,V}},
 \qquad
 \rho_{X,V,\beta}=Z_{X,V}(\beta)^{-1}e^{-\beta H_{X,V}}.
\]
El límite termodinámico existe cuando la presión
\[
 p_X(\beta)=\lim_{V\nearrow\infty}
 \frac1{\beta |V|}\log Z_{X,V}(\beta)
\]
existe y es independiente de la sucesión cofinal.

\subsection{Criterio de condensación por saturación de la densidad excitada}

La condensación bosónica se
decide por la saturación de la densidad de estados excitados; no se proclama a
partir de la mera existencia de una completación simétrica. Así, el problema
queda expresado mediante una condición matemática verificable para el
Hamiltoniano y el límite, en lugar de una atribución nominal.

\section{Sectores topológicos, excitaciones virtuales y clasificación de sus codominios}

\subsection{Categoría trenzada, datos \texorpdfstring{\(N\)--\(F\)--\(R\)}{N--F--R} y ecuaciones de pentágono y hexágono}

En un sector topológico el codominio primario es una categoría trenzada, no
una recta de masas. Para objetos simples \(a,b,c\), los datos físicos incluyen
\[
 N_{ab}^{c},\qquad F^{abc}_{d},\qquad R^{ab}_{c},
\]
con ecuaciones pentagonal y hexagonal. Una masa escalar existe únicamente
cuando una realización Hamiltoniana proporciona un proyector espectral y una
brecha. Los sectores Fibonacci, Ising/Majorana y \(\mathbb Z_6\) quedan así
cerrados como codominios categóricos; no se les asigna una cifra universal
por analogía con una partícula elemental.

En el sector \(\mathbb Z_6\), la fusión y el trenzado son
\[
 a\otimes b=a+b\pmod 6,
 \qquad
 R_{ab}=\zeta_6^{ab},
 \qquad \zeta_6=e^{2\pi i/6}.
\]
Para Fibonacci,
\[
 \tau\otimes\tau=\mathbf 1\oplus\tau,
 \qquad
 F_\tau^{\tau\tau\tau}=
 \begin{pmatrix}
 \varphi^{-1}&\varphi^{-1/2}\\
 \varphi^{-1/2}&-\varphi^{-1}
 \end{pmatrix},
\]
\[
 R_{\mathbf1}^{\tau\tau}=e^{-4\pi i/5},
 \qquad
 R_\tau^{\tau\tau}=e^{3\pi i/5}.
\]
El sector de Ising posee objetos \(\mathbf1,\psi,\sigma\) y reglas
\[
 \psi\otimes\psi=\mathbf1,
 \qquad
 \psi\otimes\sigma=\sigma,
 \qquad
 \sigma\otimes\sigma=\mathbf1\oplus\psi,
\]
\[
 F_\sigma^{\sigma\sigma\sigma}
 =\frac1{\sqrt2}
 \begin{pmatrix}1&1\\1&-1\end{pmatrix},
 \qquad
 R_{\mathbf1}^{\sigma\sigma}=e^{-\pi i/8},
 \qquad
 R_\psi^{\sigma\sigma}=e^{3\pi i/8}.
\]
Estas matrices satisfacen las relaciones de coherencia en las convenciones
indicadas. La duplicación
\(\mathcal C\boxtimes\mathcal C^{\rm rev}\) conserva ambas orientaciones y
evita introducir una quiralidad por elección externa.

\subsection{Persistencia finita de los estados virtuales y ausencia de una masa escalar universal}

Un estado virtual es una sección de persistencia finita en una amplitud o un
resolvente. No pertenece al espectro asintótico y no exige una masa de polo
real. Sus parámetros son la virtualidad, el canal y la duración de
persistencia. Esta definición evita convertir una línea interna de cálculo en
una nueva especie.

\section{Gravitación, torsión y criterio espectral del modo colectivo de espín 2}

\subsection{Exclusión taquiónica por positividad y separación de la gravitación primaria por conexión y torsión}

La positividad excluye una especie elemental con \(m^2<0\) dentro del
codominio másico. Un autovalor negativo del Hessiano de un fondo no define un
taquión como partícula: diagnostica la inestabilidad del fondo y exige cambiar
la solución alrededor de la cual se linealiza.

La gravitación entra en Mecánica Dimensional mediante conexión, curvatura y
torsión. No requiere una partícula elemental de espín \(2\) como dato de
partida.

\subsection{Polo físico del Hessiano transversal sin traza y positividad del residuo}

Sea \(\mathcal S_{\rm grav}\) la acción geométrica y
\(\mathcal H_{\rm TT}\) el subespacio transversal y sin traza de su Hessiano
en una solución. El observable linealizado es
\[
 \mathcal K_{\rm TT}
 =P_{\rm TT}\,\delta^2\mathcal S_{\rm grav}\,P_{\rm TT}.
\]
Un modo colectivo de espín \(2\) existe como partícula asintótica únicamente
si el resolvente de \(\mathcal K_{\rm TT}\) posee un polo físico con residuo
positivo y respeta las restricciones. Si no existe tal polo, la interacción
gravitatoria permanece íntegramente en la conexión y la torsión. Así, el
gravitón no es una pieza necesaria de la ley de masas, y su eventual aparición
queda decidida por un criterio espectral, no por una casilla nominal.

\section{Equivalencia entre sección persistente, límite inverso y representación finita de partícula}

\subsection{Cociente por equivalencia de cola y compatibilidad de las familias de secciones}

La definición genealógica de partícula y su representación operatorial son
dos presentaciones del mismo objeto. Sea
\[
 \mathfrak r_K=(x_K,\gamma_K,B_K,s_K,n_K,U_K)
\]
un representante finito a profundidad \(K\), con estado, ruta, frontera,
orientación, contador y representación, y sean
\(r_{L,K}:\mathfrak R_L\to\mathfrak R_K\) los mapas de refinamiento. Una
sección persistente es una familia
\[
 X=(\mathfrak r_K)_{K\ge K_0},
 \qquad r_{L,K}(\mathfrak r_L)=\mathfrak r_K
 \quad(L\ge K\ge K_0).
\]
Dos familias son equivalentes si coinciden en una cola cofinal.

\subsection{Condición de prolongación única para reconstruir una sección desde una evaluación finita}

\begin{theorem}[Equivalencia de la sección persistente y el límite inverso]
La categoría de secciones persistentes, definidas como familias compatibles,
es canónicamente equivalente a
\[
 \varprojlim_K\mathfrak R_K/\!\sim_{\rm cola}.
\]
La proyección a una profundidad finita es una evaluación, no un funtor inverso.
Un representante \(\mathfrak r_K\) determina una única sección persistente si
y sólo si posee una única prolongación compatible módulo equivalencia de cola;
esta unicidad se obtiene, por ejemplo, cuando los mapas de refinamiento son
biyectivos sobre la cola superviviente a partir de \(K\).
\end{theorem}

\begin{proof}
Por definición, una sección persistente es precisamente un punto del límite
inverso; el cociente identifica las familias que coinciden en una cola cofinal.
La evaluación es la proyección canónica a una componente. Su fibra contiene
todas las prolongaciones compatibles del representante evaluado, de modo que
es unitaria exactamente bajo la condición de prolongación única indicada.
\end{proof}

Una partícula HMT--MD es, por tanto, una clase de secciones persistentes sin
obstrucción terminal finita, provista de representación, orientación, carga,
espín y carácter espectral compatibles con el refinamiento. La ruta es su
historia observable, no su identidad reducida a una celda. La antipartícula
es la sección conjugada por la involución dodecafásica y la representación
conjugada. Una sección finita obstruida es virtual; un producto cocíclico de
secciones es compuesto; una representación trenzada persistente es anyónica.

\section{Teorema de clausura de la realización equivariante y clasificación exhaustiva de las especies materiales}

\subsection{Atlas exhaustivo de 324 rutas: 49 invariantes por sección y conservación del sector nulo}

\begin{theorem}[Cierre por codominio físico]
Supónganse construidos el atlas \(81/56/13/324\), el sector nulo, la escala
acción--reloj--masa, los lectores \(K_D,K_\Omega\), el descriptor físico y la
representación interna de cada estado. Entonces todo descriptor admisible y
realizable, es decir, todo \(X\) con
\(d_{\mathcal P}(X)\in\operatorname{im}d_{\mathcal R}\),
determina, sin consultar su valor externo:
\begin{enumerate}
\item una fibra, órbita o multisección \(\mathscr S(X)\);
\item el par positivo bidireccional y su PVM conjunta
      \((\widehat M_X^{\beta,D},\widehat M_X^{\beta,\Omega},
      E_X^{D,\Omega})\);
\item una descomposición en bloques irreducibles;
\item una medida espectral para cada acoplamiento físico;
\item un escalar en un sector comprimido de rango uno o un multiplete en un
      sector reducible; en un sector abierto, un operador efectivo y, sólo si
      existe continuación meromorfa con un cero aislado admisible, un polo
      complejo;
\item un objeto categórico cuando la realización es topológica y no posee aún
      una brecha Hamiltoniana;
\item un criterio de inexistencia cuando el tipo solicitado contradice
      positividad, calibre o asintoticidad.
\end{enumerate}
Ninguno de estos resultados requiere seleccionar una ruta mediante la masa
que después será contrastada.
\end{theorem}

\begin{proof}
El producto fibrado construye la fibra de forma prospectiva. El carácter y los
dos lectores construyen el par positivo y su PVM conjunta. Cuando el tipo
físico declara un funcional escalar, el cálculo funcional produce el operador
correspondiente; la esperanza condicional lo proyecta al conmutante de la
simetría. La descomposición isótípica produce escalares o bloques por el lema
de Schur. Los proyectores de preparación producen medidas espectrales. La
reducción de Feshbach produce el operador efectivo; la continuación meromorfa
y un cero aislado admisible producen un polo, si existe. Los sectores
topológicos conservan sus datos de fusión y trenzado hasta que exista una
realización Hamiltoniana. La positividad y las restricciones de calibre
excluyen los tipos incompatibles. Cada paso depende sólo de datos internos o
del acoplamiento declarado, nunca de la columna externa.
\end{proof}

\paragraph{Controles de refutación.}

La construcción queda refutada en su dominio si ocurre cualquiera de los
hechos siguientes:

\begin{enumerate}
\item al modificar los valores externos cambia la fibra \(\mathscr S(X)\);
\item la asignación deja de ser equivariante bajo inversión celular,
      conjugación o calibre;
\item se selecciona una celda puntual en una multisección no central sin un
      acoplamiento que rompa la simetría;
\item la síntesis de \(K_D\) y \(K_\Omega\) privilegia una orientación o no
      recupera la diagonal cuando ambos lectores coinciden;
\item una compresión con varios autovalores se reemplaza por una media sin
      proyector físico;
\item el cociente de color produce masas diferentes para rojo, verde y azul;
\item se comparan masas quark sin declarar el mismo esquema y la misma escala;
\item una fase de rango uno se presenta como mezcla PMNS después de fallar el
      contraste angular de rango completo;
\item la neutralidad se usa por sí sola para afirmar el carácter Majorana;
\item la masa de un compuesto se obtiene sumando masas de constituyentes sin
      registro de enlace;
\item la anchura se inserta como coordenada de la firma en vez de extraerse de
      un modo abierto o un polo;
\item un autovalor negativo del Hessiano se publica como partícula taquiónica;
\item la gravitación se hace depender de un gravitón elemental introducido
      antes de la conexión, la curvatura y la torsión.
\end{enumerate}

\subsection{Inventarios completos de 471 masas y 384 anchuras y reconocimiento metrológico posterior}

La comparación externa se abre únicamente después de congelar descriptor,
fibra, operador, proyector, esquema y escala. Cada fila debe distinguir masa
estable, masa de polo, estimador de resonancia y anchura. El atlas de \(324\)
fichas y los inventarios de \(471\) masas y \(384\) anchuras constituyen la
exposición exhaustiva subsiguiente. La construcción conceptual y cada ficha
se remiten recíprocamente mediante sus identificadores de ruta. Las fichas
conservan todos sus campos y se componen como unidades coherentes; no se
sustituyen por tablas históricas abreviadas.

La comparación no consiste en ampliar una lista de nombres, sino en aplicar,
para cada tipo físico, el mapa completo que conduce desde la genealogía de ruta
hasta su codominio observable.
